\documentclass[11pt]{article}
\usepackage[a4paper,margin=25mm]{geometry}
\usepackage{amsmath,amssymb,amsthm}
\usepackage[hidelinks]{hyperref}
\emergencystretch=2em
\newtheorem{theorem}{Theorem}
\title{An unavoidable even-power term in Volterra singular values\newline
\large Disproof of Conjecture 00000009778}
\author{gaochengzhecpu}\date{}
\begin{document}\maketitle
\begin{abstract}
The classical Volterra operator on $L^2(0,1)$ has ordered singular
values $s_n=1/[\pi(n-1/2)]$, $n\geq1$. Consequently
$n^2(s_n-1/(\pi n))\to1/(2\pi)\ne0$, excluding an expansion
with only odd inverse powers. A full unrestricted expansion does
exist; its even-power terms cannot be omitted. We give the
actual operator-to-spectrum argument and explicitly separate it
from the Lean-checked integral identities and asymptotic core.
The Hilbert-space spectral identification is a stated mathematical
bridge, not a claimed complete Lean formalization.
\end{abstract}

\section*{The actual compact operator}
On the real or complex Hilbert space $H=L^2(0,1)$, set
\[
 (Vf)(x)=\int_0^x f(t)\,dt.
\]
Its kernel $\mathbf1_{\{0\leq t\leq x\leq1\}}$ is precisely an
interval-indicator integration kernel, as required by the source.
Its squared $L^2$ kernel norm is $1/2$, so $V$ is a bounded
Hilbert--Schmidt operator and is compact. Every $L^2$ function on
this finite interval is integrable, and the formula selects an
absolutely continuous representative. Kernel diagonal conventions
have no effect on the operator on $L^2$ equivalence classes.

Fubini applied to the inner product gives the actual adjoint
\[
 (V^*g)(x)=\int_x^1 g(t)\,dt.
\]
Thus $A=V^*V$ is compact, positive and self-adjoint, with
\[
 (Af)(x)=\int_x^1\int_0^t f(s)\,ds\,dt
       =\int_0^1\bigl(1-\max(x,s)\bigr)f(s)\,ds.
\]
Fubini is justified, for example, by integrability of
$|f(s)||g(t)|$ on the square, using $L^2\subset L^1$ and
Cauchy--Schwarz. Positivity follows from
$\langle Af,f\rangle=\|Vf\|_2^2$.

For arbitrary $f\in H$, the displayed representative $y=Af$
satisfies
\[
 y\in C^1[0,1],\qquad y'(x)=-\int_0^x f(s)\,ds,
 \qquad y''=-f\ \hbox{a.e.},\qquad y'(0)=0,\ y(1)=0.
\]
Here $y'$ is absolutely continuous; no classical second
derivative of an arbitrary $L^2$ function is assumed.

\newpage
\section*{Complete spectral identification and ordering}
Suppose $Af=\lambda f$ in $L^2$, with $f\ne0$ and
$\lambda\ne0$. Positivity gives $\lambda>0$. Replace $f$ by
the representative $\widetilde f=y/\lambda$. It is continuous,
so the integral formula for $y'$ and the fundamental theorem
of calculus upgrade $y$ to $C^2$. Consequently
\[
 \widetilde f''+\lambda^{-1}\widetilde f=0,
 \qquad \widetilde f'(0)=0,\quad \widetilde f(1)=0
\]
hold classically everywhere. Writing $\omega=1/\sqrt\lambda>0$,
the nonzero solutions are multiples of $\cos(\omega x)$, and
$\cos\omega=0$. Therefore
\[
 \omega_n=(n-\tfrac12)\pi,\qquad
 \lambda_n=\omega_n^{-2},\qquad n=1,2,\ldots .
\]
Conversely direct integration proves
\[
 V\cos(\omega_n x)=\frac{\sin(\omega_n x)}{\omega_n},\qquad
 V^*\sin(\omega_n x)=\frac{\cos(\omega_n x)}{\omega_n}.
\]
Each listed eigenvalue thus occurs, with a one-dimensional
eigenspace by the ODE classification. There are no other nonzero
eigenvalues. The compact self-adjoint spectral theorem applies
to $A$, not to the non-self-adjoint operator $V$. It gives
exhaustiveness of the eigenvalue list with multiplicities.
Also $\ker A=\ker V=\{0\}$: if $Vf=0$ a.e., its continuous
primitive is zero everywhere, hence its derivative $f$ vanishes
a.e. The positive eigenvalues above strictly decrease to zero.
Taking their positive square roots gives the ordered singular values
\[
 s_n(V)=\sqrt{\lambda_n}=
 \frac1{\pi(n-1/2)},\qquad n\geq1.
\]
This agrees with the squared Gram spectrum stated in
Opmeer~\cite{Opmeer}, Example 31. The positive square root is
essential; the eigenvalues of $V^*V$ are not the singular values of $V$.

\section*{The missing even term}
Exact algebra yields
\[
 n^2\left(s_n-\frac1{\pi n}\right)
 =\frac{n}{\pi(2n-1)}\longrightarrow\frac1{2\pi}\ne0.
\]
An odd-only Poincar\'e expansion in $n^{-1}$ would, already
through order three, give a real $c$ such that
\[
 n^3\left(s_n-\frac1{\pi n}-\frac c{n^3}\right)\longrightarrow0.
\]
Multiplying this limit by $1/n$ and adding $c/n$ would force
the previous $n^2$-scaled difference to tend to zero, a contradiction.

\newpage
\begin{theorem}
The Volterra singular values have a nonzero $n^{-2}$ coefficient,
so they do not admit the conjectured odd-only inverse-power
expansion. This remains true under every integer change of starting index.
\end{theorem}
\begin{proof}
The first assertion follows from the nonzero limit on page 2.
For an integer shift $k$, eventually
$s_{n+k}=1/[\pi(n+a)]$ with $a=k-1/2\ne0$. Then
\[
 n^2\left(s_{n+k}-\frac1{\pi n}\right)
 =-\frac a\pi\frac n{n+a}\longrightarrow-\frac a\pi\ne0.
\]
The same necessary odd-expansion condition is impossible.
\end{proof}

The ordinary full expansion actually is geometric:
\[
 s_n=\frac1{\pi n}\left(1+\frac1{2n}+\frac1{4n^2}+\cdots\right).
\]
Every displayed coefficient is nonzero. If ``odd corrections''
were instead intended inside the relative factor, the relative
$n^{-2}$ term $1/4$ also disproves that reading. A half-integer
frequency variable $n-1/2$ is not an integer reindexing or an
expansion in the stated integer $n$. We refute the parity
restriction, not the existence of an unrestricted full expansion;
the additional Bernoulli-coefficient clause need not be settled.

\section*{Formalization boundary and reproduction}
Lean defines genuine forward and backward interval integrals on
real functions. It proves their displayed cosine-mode and Gram-mode
identities for every mode and every real argument, including the
half-integer boundary condition. Separately it proves the
nonnegative square-root step from the published Gram formula,
the actual scaled-difference limit, and failure of the necessary
third-order odd expansion. Theorems quantify over any sequence
eventually satisfying the exact spectral formula. They include
every integer offset; $n=0$ is immaterial to the eventual statements.

The actual $L^2$ bounded operator, adjoint theorem, regularity
upgrade, compact spectral theorem and exhaustive ordered spectral
identification on pages 1--2 are mathematical analysis outside
this Lean project. They are substantive bridges, not definitions
of a chosen sequence as singular values. The project does not
claim that pointwise integral identities alone prove completeness
or that the entire $L^2$ spectrum has been formalized.

\newpage
\section*{Reproduction and references}
Lean 4.19.0 is pinned, with Mathlib revision
\begin{center}
\texttt{c44e0c8ee63ca166450922a373c7409c5d26b00b}.
\end{center}
From \texttt{lean/}, run:
\begin{verbatim}
lake exe cache get
lake build
lake env lean -DwarningAsError=true Main.lean
\end{verbatim}
Run \texttt{python verify.py} at the root for supplementary exact
algebra checks. Compile with \texttt{tectonic main.tex}.
Actual fresh-build logs and standard-foundation axiom reports
are included. Author, parent and internal cross-review were
performed; no external independent review is claimed.

\begin{thebibliography}{9}
\bibitem{source} The Last Math Competition,
Conjecture 00000009778. Exact bilingual source included.
\bibitem{Opmeer} M. R. Opmeer,
\emph{Decay of singular values for infinite-dimensional systems
with Gevrey regularity}, Systems \& Control Letters 137 (2020),
104644, Example 31.
\href{https://doi.org/10.1016/j.sysconle.2020.104644}{Publisher record};
\href{https://purehost.bath.ac.uk/ws/portalfiles/portal/203458869/GevreySingularValues.pdf}{author manuscript},
PDF page 23.
\end{thebibliography}
\end{document}
